\subsection{De Sitter vs Quintessence in String Theory --- arXiv:1808.08967}

\textbf{Authors:} Michele Cicoli, Senarath de Alwis, Anshuman Maharana, Francesco Muia, Fernando Quevedo

\paragraph{主な主張}
弦理論のde Sitter構成とquintessence構成を比較し、前者の計算上の制御は改善してきた一方、後者にはより厳しい微調整や現象論的制約が課されると論じる\cite{Cicoli:2018kdo}。

\paragraph{新規性}
de Sitter構成への批判だけでなく、代替となるquintessenceにも同じ基準を適用し、Higgsポテンシャル・第五の力・超軽量アクシオンの問題を比較する。

\paragraph{位置付け}
de Sitter予想の是非を、具体的な暗黒エネルギー模型の構成可能性から検討する比較研究。

\paragraph{コメント（読書メモ）}
type IIBを含むde Sitter構成では超対称性が破れており、計算を十分に制御する手法が不足している点が問題となる。超対称性は有効理論の構造を制限するが、$N=1$でも一般の背景における量子補正の正確な形が既知とは限らない。また、超対称真空の平坦方向は、現象論に必要な超対称性の破れを導入した際にrunawayへ変わり得るため、そのまま適切な真空の出発点になるとは限らない。
